\documentclass[12pt,a4paper]{article}

% ---------- Kódovanie a jazyk ----------
\usepackage[utf8]{inputenc}
\usepackage[T1]{fontenc}
\usepackage[slovak]{babel}
\usepackage{lmodern}
\usepackage{keyval}
\usepackage[patches={}]{microtype}

% ---------- Rozloženie strany ----------
\usepackage[margin=2.5cm]{geometry}
\usepackage{setspace}
\onehalfspacing
\usepackage{parskip}
\setlength{\parindent}{0pt}
\setlength{\parskip}{0.6em}

% ---------- Zoznamy, tabuľky, odkazy ----------
\usepackage{enumitem}
\setlist{itemsep=2pt, topsep=4pt}
\usepackage{booktabs}
\usepackage{array}
\usepackage{tabularx}
\usepackage{caption}
\captionsetup{font=small, labelfont=bf, skip=6pt}
\usepackage{titlesec}
\titleformat{\section}
  {\Large\bfseries}{\thesection}{0.6em}{}
\titleformat{\subsection}
  {\large\bfseries}{\thesubsection}{0.6em}{}
\usepackage[hidelinks]{hyperref}
\usepackage{url}
\urlstyle{same}

\begin{document}

% TITULNÁ STRANA

\begin{titlepage}
    \thispagestyle{empty}

    % Univerzita a fakulta
    \begin{center}
        \vspace*{1cm}

        {\normalsize Technická univerzita v Košiciach\par}
        \vspace{0.2cm}
        {\normalsize Fakulta elektrotechniky a informatiky\par}
        \vspace{0.2cm}
        {\normalsize Predmet: Spracovanie prirodzeneho jazyka\par}

        % Názov práce
        \vspace{4cm}

        {\LARGE\bfseries
        Odpovedanie na otázky nad tabuľkovými dátami\par}

        \vspace{0.4cm}

        {\Large s využitím datasetu DataBench\par}

        \vspace{0.5cm}

    \end{center}

    % Autori a študijný program
    \vfill

    \begin{flushleft}
        \textbf{Autori:}\\
        Volodymyr Ihnatov\\
        Artem Marynov\\
        Ruslan Savka\\
        Kostiantyn Ostrovskyi

        \vspace{0.5cm}

        \textbf{Študijný program:}\\
        Umelá inteligencia
    \end{flushleft}

    % Miesto a rok
    \vspace{2cm}

    \begin{center}
        Košice, 2026
    \end{center}

\end{titlepage}


% OBSAH

\tableofcontents
\newpage

% JEDNOTLIVÉ KAPITOLY

\section{Úvod a cieľ projektu}

V súčasnosti zohráva spracovanie prirodzeného jazyka (Natural Language Processing, NLP) dôležitú úlohu pri vývoji inteligentných systémov, ktoré dokážu porozumieť používateľským požiadavkám a poskytovať relevantné odpovede. Jednou z aktuálnych úloh v tejto oblasti je automatické odpovedanie na otázky o tabuľkových dátach formulované v prirodzenom jazyku. Riešenie takýchto úloh si vyžaduje nielen porozumenie textovej otázke, ale aj schopnosť analyzovať štruktúru tabuľky, vyhľadávať potrebné informácie a vykonávať výpočty.

Cieľom tohto projektu je preskúmať možnosti jazykových modelov pri riešení úloh automatického odpovedania na otázky nad tabuľkovými dátami s využitím datasetu DataBench. Vstupom systému je otázka formulovaná v prirodzenom jazyku a príslušná tabuľka. Výstupom je odpoveď vo forme logickej hodnoty (Boolean), kategórie, čísla alebo zoznamu hodnôt. Hlavná pozornosť sa venuje presnosti odpovedí a schopnosti systému správne interpretovať otázky a získavať potrebné informácie z tabuliek.

Hlavnými cieľmi projektu sú:
\begin{enumerate}
    \item analyzovať úlohu DataBench, jej štruktúru a typy otázok a odpovedí,
    \item navrhnúť a implementovať základné riešenie na automatické generovanie odpovedí na otázky týkajúce sa tabuľkových dát,
    \item určiť vhodné metriky hodnotenia a overiť kvalitu základného riešenia na dostupných dátach,
    \item analyzovať výhody a obmedzenia zvoleného prístupu a identifikovať možnosti jeho zlepšenia v ďalšej etape projektu.
\end{enumerate}

\textbf{Výskumná otázka projektu:} Do akej miery dokáže jazykový model efektívne odpovedať na otázky o tabuľkových dátach a aké prístupy môžu zvýšiť presnosť získaných odpovedí?

\section{Analýza problému a datasetu DataBench}

\subsection{Analýza riešeného problému}

Projekt sa zaoberá automatickým odpovedaním na otázky formulované v prirodzenom jazyku na základe štruktúrovaných tabuľkových dát. Táto úloha patrí do oblasti spracovania prirodzeného jazyka (NLP) a vyžaduje kombináciu porozumenia textu, analýzy dát a logického uvažovania \cite{semeval}.

Vstupom systému je otázka v prirodzenom jazyku a tabuľka obsahujúca relevantné údaje. Na základe otázky musí systém identifikovať potrebné stĺpce a riadky, správne interpretovať požadovanú operáciu a vytvoriť odpoveď, ktorá zodpovedá obsahu tabuľky. V závislosti od otázky môže byť potrebné vyhľadať konkrétnu hodnotu, porovnať údaje, vypočítať štatistiku alebo vytvoriť zoznam výsledkov.

Hlavnou výzvou je zabezpečiť, aby systém správne porozumel otázke a zároveň vykonal potrebné operácie nad dátami. Nestačí teda iba generovať jazykovo správnu odpoveď. Odpoveď musí byť podložená informáciami z poskytnutej tabuľky. Systém by nemal dopĺňať informácie z externých zdrojov ani vytvárať hodnoty, ktoré z dostupných dát nevyplývajú \cite{semeval}.

\subsection{Charakteristika datasetu DataBench}

Na riešenie uvedeného problému sa používa dataset DataBench, ktorý je určený na hodnotenie schopnosti systémov odpovedať na otázky týkajúce sa tabuľkových dát \cite{semeval,hfdatabench}. Dataset poskytuje otázky a identifikátory príslušných datasetov, na základe ktorých možno určiť, s ktorou tabuľkou sa má pri odpovedaní pracovať.

Testovacie dáta sú dostupné v archíve \texttt{competition\_databench.zip}. Archív obsahuje súbor \texttt{test\_qa.csv}, v ktorom sú uložené otázky a identifikátory zodpovedajúcich datasetov. Každému identifikátoru datasetu zodpovedá priečinok obsahujúci dva súbory:
\begin{enumerate}
    \item \texttt{all.parquet} -- kompletná verzia tabuľkových dát určená pre DataBench,
    \item \texttt{sample.parquet} -- vzorka obsahujúca prvých 20 riadkov zodpovedajúcej tabuľky, určená pre DataBench Lite \cite{hfdatabench}.
\end{enumerate}

Dataset teda umožňuje skúmať odpovedanie na otázky nad kompletnými tabuľkami aj nad ich menšími vzorkami. Takéto rozdelenie poskytuje možnosť porovnať úspešnosť systému pri práci s rôznym množstvom dostupných dát.

Okrem testovacích dát sú k dispozícii aj vývojové a tréningové dáta publikované na platforme Hugging Face v datasete \texttt{cardiffnlp/databench}. Tieto dáta možno využiť pri príprave a testovaní riešenia pred jeho konečným vyhodnotením \cite{hfdatabench}.

\subsection{Typy otázok a očakávaných odpovedí}

Podľa dokumentácie sú odpovede v DataBench rozdelené do piatich základných typov:
\begin{enumerate}
    \item \textbf{Boolean} -- odpoveď predstavuje logickú hodnotu, napríklad \texttt{True} alebo \texttt{False}. Akceptované sú aj ekvivalentné formy Yes/No a Y/N.
    \item \textbf{Category} -- odpoveď tvorí kategória alebo hodnota získaná z bunky tabuľky, prípadne podreťazec jej obsahu.
    \item \textbf{Number} -- odpoveď je číselná hodnota získaná z tabuľky alebo vypočítaná na základe dostupných údajov, napríklad pri výpočte priemeru, maxima alebo minima.
    \item \textbf{List[category]} -- odpoveď obsahuje zoznam kategórií. Pri tvorbe odpovede je potrebné zohľadniť, či otázka vyžaduje jedinečné hodnoty alebo povoľuje opakovanie.
    \item \textbf{List[number]} -- odpoveď obsahuje zoznam číselných hodnôt získaných z tabuľky alebo výpočtom.
\end{enumerate}

Správne určenie typu očakávanej odpovede je dôležitou súčasťou riešenia. Pomáha zvoliť vhodný postup spracovania otázky a zabezpečiť, aby výsledok mal požadovaný formát. Pri vyhodnocovaní sa poradie prvkov v zoznamových odpovediach nepovažuje za rozhodujúce \cite{hfdatabench,databencheval}.

\subsection{Obmedzenia a spôsob vyhodnocovania}

Pri práci s datasetom je potrebné dodržiavať pravidlo, podľa ktorého sa odpovede musia zakladať výlučne na poskytnutých tabuľkových dátach. Používanie externých dátových zdrojov na získanie odpovedí nie je povolené. Toto obmedzenie umožňuje hodnotiť schopnosť systému pracovať priamo s dostupnými údajmi.

Výsledky sa vyhodnocujú pomocou automatizovaného hodnotiaceho nástroja založeného na balíku \texttt{databench\_eval}. Predikcie systému sa porovnávajú s referenčnými odpoveďami. Hodnotiaca funkcia akceptuje aj niektoré odlišnosti vo formátovaní, aby sa znížil vplyv drobných formálnych chýb na výsledné hodnotenie \cite{hfdatabench}.

Výstup riešenia musí obsahovať dva súbory: \texttt{predictions.txt} pre kompletné tabuľky DataBench a \texttt{predictions\_lite.txt} pre vzorky DataBench Lite. Každý riadok obsahuje jednu odpoveď v rovnakom poradí, v akom sú otázky uvedené v testovacom súbore \cite{hfdatabench,databencheval}.

Z pohľadu projektu bude dôležité sledovať presnosť odpovedí a analyzovať typy chýb, ktoré vznikajú pri interpretácii otázok, výbere údajov alebo vykonávaní výpočtov. Tieto zistenia budú slúžiť ako východisko pre návrh a implementáciu základného riešenia (baseline) a pre následné experimenty v druhej etape projektu.

\subsection{Podrobná štruktúra datasetu}

Cieľom datasetu DataBench je umožniť systematické hodnotenie jazykových modelov pri odpovedaní na otázky nad tabuľkovými dátami. Obsahuje tabuľky z rôznych tematických oblastí a otázky vytvorené na základe ich obsahu \cite{semeval}.

Aktuálna verzia datasetu na platforme Hugging Face obsahuje 80 tabuliek a 1\,822 otázok. Celkový počet riadkov v zdrojových tabuľkách je 3\,708\,884 a počet stĺpcov je 2\,006 \cite{hfdatabench}. Tieto údaje sa týkajú aktuálnej verzie repozitára a nemali by sa zamieňať s pôvodnou verziou benchmarku z roku 2024, ktorá obsahovala 65 tabuliek a 1\,300 otázok.

\begin{table}[htbp]
\centering
\caption{Základné charakteristiky datasetu DataBench}
\label{tab:basic}
\begin{tabularx}{\textwidth}{@{}>{\raggedright\arraybackslash}p{6cm} X@{}}
\toprule
\textbf{Charakteristika} & \textbf{Popis} \\
\midrule
Názov datasetu & DataBench \\
Hlavná úloha & Odpovedanie na otázky nad tabuľkovými dátami \\
Počet tabuliek v aktuálnej verzii & 80 \\
Počet otázok v aktuálnej verzii & 1\,822 \\
Celkový počet riadkov zdrojových tabuliek & 3\,708\,884 \\
Celkový počet stĺpcov zdrojových tabuliek & 2\,006 \\
Formát dát & Parquet \\
Jazyky uvedené v repozitári & Angličtina a španielčina \\
Hlavné verzie tabuliek & Kompletné tabuľky (\texttt{all.parquet}) a vzorky (\texttt{sample.parquet}) \\
Zdroj & Hugging Face -- \texttt{cardiffnlp/databench} \\
\bottomrule
\end{tabularx}
\end{table}

Každý záznam v časti datasetu určenej na otázky a odpovede obsahuje informácie potrebné na identifikáciu otázky, príslušnej tabuľky a správnej odpovede. V aktuálnej verzii repozitára sú v náhľade dostupné stĺpce uvedené v tabuľke~\ref{tab:structure} \cite{hfdatabench}.

\begin{table}[htbp]
\centering
\caption{Štruktúra jednotlivých záznamov}
\label{tab:structure}
\begin{tabularx}{\textwidth}{@{}>{\ttfamily\raggedright\arraybackslash}p{4cm} X@{}}
\toprule
\textbf{Stĺpec} & \textbf{Význam} \\
\midrule
question & Otázka formulovaná v prirodzenom jazyku \\
answer & Referenčná odpoveď pre kompletnú tabuľku \\
type & Typ očakávanej odpovede \\
columns\_used & Názvy stĺpcov potrebných na zodpovedanie otázky \\
column\_types & Dátové typy použitých stĺpcov \\
sample\_answer & Referenčná odpoveď pre DataBench Lite \\
dataset & Identifikátor tabuľky, napríklad \texttt{001\_Forbes} \\
\bottomrule
\end{tabularx}
\end{table}

Identifikátor v stĺpci \texttt{dataset} umožňuje načítať príslušnú tabuľku zo zodpovedajúceho priečinka. Stĺpec \texttt{answer} obsahuje odpoveď pre kompletnú tabuľku, zatiaľ čo \texttt{sample\_answer} predstavuje odpoveď pre vzorku obsahujúcu prvých 20 riadkov. Tieto odpovede sa môžu líšiť, pretože výpočty nad celou tabuľkou a nad jej vzorkou môžu viesť k odlišným výsledkom \cite{hfdatabench}.

Zdrojové tabuľky sú uložené vo formáte Parquet. Súbor \texttt{all.parquet} obsahuje spracovanú kompletnú tabuľku a súbor \texttt{sample.parquet} jej zmenšenú verziu. Okrem toho repozitár obsahuje súbory s otázkami a doplňujúce informácie o jednotlivých tabuľkách.

Pri práci s DataBench je potrebné rozlišovať medzi pôvodnou kolekciou otázok a rozdeleniami určenými pre súťaž SemEval 2025 Task 8. Podľa dokumentácie Hugging Face možno načítať konfiguráciu \texttt{qa} s rozdelením \texttt{train} a konfiguráciu \texttt{semeval} s rozdeleniami \texttt{train}, \texttt{dev} a \texttt{test} \cite{hfdatabench}.

\begin{table}[htbp]
\centering
\caption{Dostupné konfigurácie a rozdelenia}
\label{tab:configs}
\begin{tabularx}{\textwidth}{@{}>{\ttfamily}l >{\ttfamily}l X@{}}
\toprule
\textbf{Konfigurácia} & \textbf{Rozdelenie} & \textbf{Účel} \\
\midrule
qa      & train & Pôvodná kolekcia otázok a odpovedí \\
semeval & train & Tréningové dáta súťaže \\
semeval & dev   & Vývojové dáta na testovanie a ladenie riešenia \\
semeval & test  & Testovacie otázky určené na konečné vyhodnotenie \\
\bottomrule
\end{tabularx}
\end{table}

Počet otázok v každom rozdelení konfigurácie \texttt{semeval} je potrebné zistiť priamo načítaním príslušnej verzie datasetu. Počty 1\,300 a 1\,822 uvedené v literatúre predstavujú celkové počty otázok v rôznych verziách DataBench, nie počty otázok v jednotlivých súťažných rozdeleniach \cite{semeval,hfdatabench}.

Pre reprodukovateľnosť experimentov je preto vhodné uviesť názov konfigurácie, názov rozdelenia a skutočný počet záznamov získaný pri načítaní datasetu.

Nasledujúce príklady pochádzajú z náhľadu oficiálneho datasetu na Hugging Face. Uvádzané názvy stĺpcov predstavujú relevantnú časť tabuľky potrebnú na zodpovedanie otázky, nie úplný výpis všetkých jej riadkov.

\paragraph{Príklad 1 -- odpoveď typu Boolean}
\begin{itemize}
    \item Identifikátor tabuľky: \texttt{001\_Forbes}
    \item Otázka: \emph{Is the person with the highest net worth self-made?}
    \item Použité stĺpce: \texttt{finalWorth}, \texttt{selfMade}
    \item Typ odpovede: \texttt{boolean}
    \item Správna odpoveď pre kompletnú tabuľku (\texttt{answer}): \texttt{True}
    \item Odpoveď pre DataBench Lite (\texttt{sample\_answer}): \texttt{False}
\end{itemize}
Tento príklad ukazuje, že systém musí identifikovať osobu s najvyšším čistým majetkom a následne overiť hodnotu príslušného atribútu \texttt{selfMade}. Rozdiel medzi odpoveďami pre kompletnú tabuľku a vzorku ilustruje vplyv rozsahu dostupných dát.

\paragraph{Príklad 2 -- číselná odpoveď}
\begin{itemize}
    \item Identifikátor tabuľky: \texttt{001\_Forbes}
    \item Otázka: \emph{What is the age of the youngest billionaire?}
    \item Použitý stĺpec: \texttt{age}
    \item Typ odpovede: \texttt{number}
    \item Správna odpoveď pre kompletnú tabuľku (\texttt{answer}): \texttt{19.0}
    \item Odpoveď pre DataBench Lite (\texttt{sample\_answer}): \texttt{32.0}
\end{itemize}
Na zodpovedanie otázky je potrebné nájsť najnižšiu hodnotu veku v stĺpci \texttt{age}. Výsledok sa líši podľa toho, či systém pracuje s kompletnou tabuľkou alebo iba s prvými 20 riadkami.

\paragraph{Príklad 3 -- zoznam kategórií}
\begin{itemize}
    \item Identifikátor tabuľky: \texttt{001\_Forbes}
    \item Otázka: \emph{List the top 3 countries with the most billionaires.}
    \item Použitý stĺpec: \texttt{country}
    \item Typ odpovede: \texttt{list[category]}
    \item Správna odpoveď pre kompletnú tabuľku (\texttt{answer}): \texttt{['United States', 'China', 'India']}
    \item Odpoveď pre DataBench Lite (\texttt{sample\_answer}): \texttt{['United States', 'China', 'Brazil']}
\end{itemize}
Tento príklad vyžaduje spočítať počet miliardárov pre jednotlivé krajiny a vybrať tri krajiny s najvyšším počtom záznamov.

Uvedené príklady ilustrujú tri odlišné typy odpovedí a rôzne operácie nad tabuľkovými dátami. Zároveň ukazujú, prečo je potrebné pri vyhodnocovaní presne rozlišovať medzi kompletnou verziou DataBench a jej zmenšenou verziou DataBench Lite.

\section{Prehľad existujúcich riešení a súvisiacich prác}

\subsection{Pôvodná práca o DataBench}

Základným zdrojom pre tento projekt je práca Osés Grijalba et al. (2024) s názvom \emph{Question Answering over Tabular Data with DataBench: A Large-Scale Empirical Evaluation of LLMs} \cite{databench2024}. Autori sa zaoberajú problémom hodnotenia jazykových modelov pri odpovedaní na otázky nad tabuľkovými dátami. Keďže existujúce benchmarky neposkytovali dostatočne rozmanité a rozsiahle dáta, autori vytvorili DataBench, ktorý v pôvodnej verzii obsahuje 65 tabuliek z rôznych oblastí a 1\,300 otázok s referenčnými odpoveďami.

Autori porovnávali modely využívajúce generovanie kódu aj modely pracujúce priamo s otázkou a tabuľkou v kontexte. Výsledky ukázali, že oba prístupy majú svoje obmedzenia a že ani jednoduché otázky nemusia byť pre jazykové modely vždy triviálne. Táto práca je pre náš projekt dôležitá, pretože definuje riešenú úlohu, predstavuje použitý benchmark a poskytuje východisko na porovnávanie rôznych prístupov.

\subsection{Hodnotiaci framework \texttt{databench\_eval}}

Oficiálny repozitár \texttt{databench\_eval} poskytuje framework na experimentovanie a hodnotenie systémov určených na odpovedanie na otázky nad tabuľkovými dátami \cite{databencheval}. Jeho cieľom je zjednotiť základné časti experimentu: vytvorenie vstupného promptu z otázky a tabuľky, zavolanie modelu a vyhodnotenie vytvorenej odpovede.

Framework umožňuje používať rôzne modely a experimentovať s rôznymi spôsobmi tvorby promptov bez nutnosti implementovať celý hodnotiaci proces od začiatku. Hodnotiaca funkcia je oddelená od samotného modelu, pretože spôsob posudzovania správnosti závisí aj od typu odpovede.

Pre projekt je tento nástroj vhodný na reprodukovateľné vyhodnotenie baseline riešenia. Umožní nám porovnať predikcie s referenčnými odpoveďami a sledovať, ako sa mení úspešnosť pri úpravách spôsobu spracovania otázok.

\subsection{Priame odpovedanie nad tabuľkou: TAPAS}

Herzig et al. (2020) predstavili model TAPAS v práci \emph{TaPas: Weakly Supervised Table Parsing via Pre-training} \cite{tapas}. Tento prístup rieši odpovedanie na otázky nad tabuľkami bez potreby explicitne generovať SQL alebo inú formálnu reprezentáciu programu. Model založený na architektúre BERT vyberá relevantné bunky tabuľky a podľa potreby predikuje aj agregačnú operáciu.

Výhodou TAPAS je priame prepojenie prirodzeného jazyka so štruktúrou tabuľky. Nevýhodou v kontexte nášho projektu je potreba vhodného predtrénovaného modelu a prispôsobenie jeho výstupov rôznym typom odpovedí v DataBench. Pri zložitých otázkach, ktoré vyžadujú rozsiahlejšie výpočty alebo viacero operácií, môže byť potrebné použiť iný prístup.

TAPAS predstavuje vhodný referenčný prístup na porovnanie s riešením založeným na generovaní kódu. Pomáha tiež vysvetliť rozdiel medzi priamym výberom buniek a vykonávaním explicitne vytvorených operácií nad tabuľkou.

\subsection{Text-to-SQL a sémantická analýza}

Ďalším významným smerom je prevod otázok v prirodzenom jazyku na formálne databázové dotazy. Tento prístup sa označuje ako Text-to-SQL. Model najskôr interpretuje otázku a vytvorí SQL dotaz, ktorý sa následne vykoná nad databázou. Výsledok vykonania dotazu predstavuje odpoveď na otázku.

Významným benchmarkom pre tento smer je Spider od Yu et al. (2018), ktorý obsahuje 10\,181 otázok a 5\,693 unikátnych SQL dotazov nad 200 databázami z rôznych oblastí \cite{spider}. Na rozdiel od jednoduchších benchmarkov Spider kladie dôraz na zovšeobecňovanie na nové databázové schémy a zložitejšie dotazy.

Výhodou Text-to-SQL je možnosť presne definovať operácie, ako sú filtrovanie, zoskupovanie, agregácia a triedenie. Tento prístup je vhodný najmä vtedy, keď sú tabuľkové dáta uložené v relačnej databáze. Pri DataBench však treba zabezpečiť prevod tabuliek do vhodného databázového formátu a správne mapovanie názvov stĺpcov na SQL dotazy.

Pre náš projekt predstavuje Text-to-SQL alternatívu ku generovaniu Python kódu. Oba prístupy vytvárajú vykonateľnú reprezentáciu požadovanej operácie, ale líšia sa použitým jazykom a prostredím vykonávania.

\subsection{Generovanie Python kódu na zodpovedanie otázok}
\label{subsec:python-code-generation}

Ďalším prístupom k riešeniu úloh typu Table Question Answering (Table QA) je generovanie zdrojového kódu na základe otázky formulovanej v prirodzenom jazyku. Jazykový model namiesto priameho vytvorenia odpovede vygeneruje napríklad Python kód, ktorý následne spracuje tabuľkové údaje a vypočíta požadovaný výsledok. Takýto prístup umožňuje využívať operácie knižníc na spracovanie dát, ako sú filtrovanie riadkov, triedenie, agregácia hodnôt a výber jedinečných kategórií.

Pri práci s tabuľkami uloženými vo formáte Parquet môže vygenerovaný kód využívať knižnicu \texttt{pandas} na načítanie údajov a ich ďalšie spracovanie. Výhodou tohto riešenia je, že samotný jazykový model nemusí vykonávať všetky výpočty priamo v texte odpovede. Namiesto toho vytvorí postup, ktorého výsledok možno vypočítať pomocou programového kódu. Príbuzný prístup využíva aj projekt \textit{Question Answering over Tabular Data with LLM Code Generation}, ktorý obsahuje implementáciu generovania kódu pre úlohy nad tabuľkovými dátami a podporuje experimenty s datasetom DataBench \cite{tabularqa}.

Tento prístup však prináša aj určité riziká. Vygenerovaný kód môže obsahovať syntaktické chyby, používať nesprávne stĺpce alebo implementovať nesprávnu interpretáciu otázky. Preto je potrebné kontrolovať výsledok vykonania kódu, ošetriť prípadné výnimky a obmedziť prostredie, v ktorom sa kód vykonáva. Pri experimentoch je zároveň dôležité zaznamenávať chyby, aby bolo možné analyzovať príčiny nesprávnych odpovedí.

Pre tento projekt je generovanie Python kódu vhodným prístupom, pretože umožňuje pracovať priamo s pôvodnými tabuľkovými údajmi a riešiť otázky vyžadujúce výpočty, triedenie alebo filtrovanie. Zároveň poskytuje jasne definovateľný postup, ktorý možno opakovane testovať a vyhodnocovať.

\subsection{Porovnanie existujúcich prístupov}
\label{subsec:comparison-approaches}

Jednotlivé prístupy k úlohe Table QA sa líšia spôsobom, akým interpretujú otázku a získavajú odpoveď z tabuľkových údajov. Model TAPAS pracuje priamo s tabuľkou a vyberá relevantné bunky, prípadne používa agregačnú operáciu. Prístup Text-to-SQL prevádza otázku do formálneho databázového dotazu, ktorý sa následne vykoná nad databázou. Generovanie Python kódu namiesto toho umožňuje použiť všeobecné programovacie operácie na spracovanie dát. Pri priamom generovaní odpovede model vytvára výsledok bez samostatného vykonateľného programu.

\begin{table}[htbp]
    \centering
    \caption{Porovnanie prístupov k riešeniu úloh Table QA}
    \label{tab:approaches-comparison}
    \small
    \renewcommand{\arraystretch}{1.2}
    \begin{tabularx}{\textwidth}{
        >{\raggedright\arraybackslash}p{3cm}
        >{\raggedright\arraybackslash}X
        >{\raggedright\arraybackslash}X}
        \toprule
        \textbf{Prístup} &
        \textbf{Princíp} &
        \textbf{Výhody a obmedzenia} \\
        \midrule
        TAPAS &
        Výber buniek tabuľky a prípadná agregácia &
        Priama práca s tabuľkou; vyžaduje vhodný predtrénovaný
        model a môže byť obmedzený typom úlohy. \\

        Text-to-SQL &
        Prevod otázky na SQL dotaz &
        Vhodný na filtrovanie a agregácie v databázach;
        vyžaduje správne vytvorený a vykonateľný dotaz. \\

        Generovanie Python kódu &
        Vytvorenie a vykonanie programu nad tabuľkou &
        Flexibilné spracovanie dát; rizikom sú chyby
        v kóde a potreba bezpečného vykonávania. \\

        Priama odpoveď LLM &
        Vygenerovanie odpovede v prirodzenom jazyku &
        Jednoduchá implementácia; model môže nesprávne
        interpretovať údaje alebo vytvoriť nepresný výsledok. \\
        \bottomrule
    \end{tabularx}
\end{table}

Pre tento projekt je dôležité nielen dosiahnuť správnu odpoveď, ale aj zabezpečiť reprodukovateľnosť experimentu. Pri priamom generovaní odpovede môže byť ťažšie určiť, ako model k výsledku dospel. Generovanie Python kódu poskytuje medzikrok, ktorý možno zaznamenať, skontrolovať a opätovne vykonať. Zároveň umožňuje pracovať priamo s tabuľkovými dátami pomocou knižnice \texttt{pandas} a riešiť nielen jednoduché vyhľadávanie hodnôt, ale aj filtrovanie, agregácie a rôzne výpočty. V porovnaní s prístupmi založenými na Text-to-SQL alebo špecializovaných tabuľkových modeloch je tento prístup jednoduchšie prispôsobiť konkrétnemu formátu datasetu DataBench.

Nevýhodou generovania kódu je potreba bezpečne kontrolovať jeho vykonávanie a spracovávať prípadné chyby. Tento problém je v navrhovanom riešení riešený pomocou izolovaného sandboxu, v ktorom sa vygenerovaný kód vykonáva oddelene od hlavného systému. Na základe uvedených vlastností sme ako hlavný baseline zvolili prístup založený na generovaní Python kódu. Priame generovanie odpovede jazykovým modelom môže byť následne použité ako porovnávací variant.

\subsection{Navrhovaný baseline}
\label{subsec:baseline}

Ako základné riešenie (baseline) plánujeme implementovať systém, v ktorom jazykový model generuje Python kód na zodpovedanie otázok nad tabuľkovými dátami. Tento prístup je v súlade s koncepciou systému \textit{Question Answering over Tabular Data with LLM Code  Generation}, ktorý využíva prevod otázok do Python kódu a jeho vykonanie nad tabuľkovými údajmi \cite{tabularqa}.

Vstupom systému bude otázka z datasetu DataBench a identifikátor príslušnej tabuľky. Tabuľkové údaje sa načítajú zo súboru vo formáte Parquet pomocou knižnice \texttt{pandas}. Model dostane text otázky, názvy stĺpcov a podľa potreby aj informácie o dátových typoch. Na základe týchto informácií vytvorí Python kód, ktorý vypočíta požadovanú odpoveď.

Navrhovaný postup bude pozostávať z nasledujúcich krokov:

\begin{enumerate}
    \item Načítanie otázky a identifikátora tabuľky z datasetu
    DataBench.
    \item Načítanie príslušnej tabuľky zo súboru Parquet.
    \item Príprava promptu obsahujúceho otázku a schému tabuľky.
    \item Generovanie Python kódu jazykovým modelom.
    \item Kontrolované vykonanie vygenerovaného kódu nad
    tabuľkovými údajmi.
    \item Konverzia výsledku do požadovaného formátu odpovede.
    \item Vyhodnotenie predikcie pomocou nástroja
    \texttt{databench\_eval}.
\end{enumerate}

Osobitná pozornosť sa bude venovať formátu výstupu, pretože DataBench obsahuje rôzne typy odpovedí, napríklad booleovské hodnoty, kategórie, čísla a zoznamy hodnôt. Výsledok preto bude potrebné normalizovať tak, aby zodpovedal požiadavkám evaluačného nástroja. Ak sa pri generovaní alebo vykonávaní kódu vyskytne chyba, systém ju zaznamená a spracuje bez neoprávneného nahradenia výsledku správnou odpoveďou \cite{oseibrefo}.

Na vyhodnotenie sa použije oficiálny framework \texttt{databench\_eval}, ktorý poskytuje infraštruktúru na spúšťanie modelov a vyhodnocovanie ich odpovedí. Na zabezpečenie porovnateľnosti budú jednotlivé varianty testované na rovnakej množine otázok a pri rovnakých podmienkach. Ako doplnkový experiment sa plánuje porovnať generovanie Python kódu s priamym generovaním odpovedí bez vykonávania kódu.

Výsledky budú analyzované pomocou metrík podporovaných evaluačným frameworkom. Okrem celkovej úspešnosti bude užitočné zaznamenávať aj počet chýb pri generovaní a vykonávaní kódu. Takáto analýza pomôže identifikovať, či nesprávne odpovede vznikajú najmä pri interpretácii otázky, výbere stĺpcov, výpočte výsledku alebo konverzii výstupu.

Keďže ide o plánovaný baseline, jeho skutočná úspešnosť a výpočtová náročnosť budú uvedené až po implementácii a vykonaní experimentov.



\section{Implementácia navrhnutého riešenia}

V tejto kapitole je opísaná implementácia základného riešenia pre odpovedanie
na otázky nad tabuľkovými dátami z datasetu DataBench. Implementované riešenie
využíva jazykový model Gemini na generovanie Python kódu, ktorý následne
vykonáva požadovanú operáciu nad tabuľkovými dátami. Výsledok vykonania je
normalizovaný a následne porovnaný s referenčnou odpoveďou pomocou nástroja
\texttt{databench\_eval}.

Implementácia je rozdelená do viacerých samostatných modulov. Každý modul
zodpovedá za konkrétnu časť spracovania, čo umožňuje jednoduchšie testovanie
a prípadné úpravy jednotlivých častí systému.

Projekt je rozdelený do niekoľkých Python súborov. Hlavný priebeh programu
zabezpečuje súbor \texttt{main.py}, ktorý spája jednotlivé komponenty do
jedného procesu. Načítanie tabuliek zabezpečuje modul
\texttt{dataset\_loader.py}, tvorbu promptu modul
\texttt{prompt\_builder.py} a komunikáciu s jazykovým modelom modul
\texttt{llm\_client.py}.

Vykonávanie vygenerovaného kódu je oddelené do modulov
\texttt{code\_executor.py} a \texttt{sandbox\_runner.py}. Modul
\texttt{evaluator.py} zabezpečuje porovnanie výslednej odpovede s referenčnou
odpoveďou a modul \texttt{metrics.py} vypočítava základné metriky z uložených
výsledkov.

\begin{verbatim}
DataBench
    |
    v
main.py
    |
    +--> dataset_loader.py
    |       |
    |       +--> all.parquet / sample.parquet
    |
    +--> prompt_builder.py
    |       |
    |       +--> vytvorenie promptu
    |
    +--> llm_client.py
    |       |
    |       +--> Gemini
    |       |
    |       +--> Python kód
    |
    +--> code_executor.py
    |       |
    |       +--> kontrola syntaxe
    |       +--> Docker
    |              |
    |              +--> sandbox_runner.py
    |
    +--> evaluator.py
    |       |
    |       +--> porovnanie odpovede
    |
    +--> results.jsonl
            |
            +--> metrics.py
\end{verbatim}

Takéto rozdelenie umožňuje oddeliť načítanie dát, generovanie kódu, jeho
vykonanie a vyhodnotenie výsledkov.

\subsection{Načítanie otázok a tabuľkových dát}

Hlavný program načítava otázky z datasetu DataBench pomocou knižnice
\texttt{datasets}. V aktuálnej implementácii sa používa konfigurácia
\texttt{qa} a split \texttt{train}. Každá otázka obsahuje okrem samotného
textu otázky aj identifikátor tabuľky, typ očakávanej odpovede a referenčnú
odpoveď.

Pre každú otázku sa podľa identifikátora datasetu vyhľadá príslušná lokálna
tabuľka vo formáte Parquet. Implementácia podporuje dva režimy spracovania.
V režime \texttt{full} sa používa súbor \texttt{all.parquet}, ktorý obsahuje
celú tabuľku. V režime \texttt{lite} sa používa súbor
\texttt{sample.parquet}, ktorý predstavuje menšiu vzorku dát.

Na načítanie tabuľky sa používa knižnica \texttt{pandas}. Po načítaní tabuľky
sa pomocou funkcie \texttt{get\_table\_schema()} vytvorí schéma obsahujúca
názvy jednotlivých stĺpcov a ich dátové typy. Táto informácia sa následne
používa pri tvorbe promptu pre jazykový model.

\subsection{Tvorba promptu a Generovanie Python kódu}

Po načítaní otázky a tabuľky sa vytvorí prompt pomocou funkcie
\texttt{build\_prompt()}. Prompt obsahuje samotnú otázku, zoznam dostupných
stĺpcov, ich dátové typy a očakávaný typ odpovede.

Dôležitou súčasťou promptu sú pravidlá pre generovanie výsledného Python
kódu. Jazykový model dostáva informáciu, že objekt \texttt{df} obsahujúci
tabuľku je už pripravený. Vygenerovaný program musí uložiť výsledok do
premennej \texttt{result}.

Prompt zároveň obsahuje pravidlá pre jednotlivé typy odpovedí. Podporované
sú napríklad booleovské hodnoty, čísla, reťazce, kategórie a zoznamy
kategórií alebo čísel. Týmto spôsobom sa znižuje riziko, že jazykový model
vráti textové vysvetlenie namiesto požadovanej hodnoty.

V prompte sú tiež uvedené pravidlá týkajúce sa spracovania duplicitných
hodnôt, zoradenia výsledkov, číselných hodnôt a konverzie dátových typov.
Model je zároveň upozornený, že nesmie používať stĺpce alebo hodnoty, ktoré
nie sú súčasťou poskytnutej tabuľky.


Na generovanie Python kódu sa používa model Gemini prostredníctvom knižnice
\texttt{google-genai}. API kľúč je načítaný z environmentálnej premennej
\texttt{GEMINI\_API\_KEY}, ktorá je uložená v súbore \texttt{.env}.

Model dostane pripravený prompt a jeho úlohou je vytvoriť vykonateľný Python
kód. Výstup modelu je následne upravený pomocou funkcie
\texttt{clean\_generated\_code()}, ktorá odstraňuje prípadné Markdown
ohraničenie kódu, napríklad bloky označené pomocou \texttt{```python}.

Implementácia obsahuje aj opakovanie požiadavky pri dočasných chybách
komunikácie s API. Pri chybách súvisiacich s prekročením limitu požiadaviek
sa používa postupné predlžovanie času čakania medzi jednotlivými pokusmi.
Maximálny počet pokusov komunikácie s modelom je nastavený na päť.

Pred samotným vykonaním kódu sa kontroluje jeho syntaktická správnosť.
Na tento účel sa používa modul \texttt{ast} a funkcia
\texttt{ast.parse()}.

Ak kód obsahuje syntaktickú chybu, nie je odoslaný na vykonanie. Systém
zaznamená chybu a môže požiadať jazykový model o novú verziu kódu. V rámci
spracovania jednej otázky sú povolené maximálne dve generovania kódu --
prvotný pokus a jedna opakovaná generácia.

Pri opakovanom pokuse je model informovaný o tom, že predchádzajúci
vygenerovaný kód nebol úspešný. Do promptu sa vloží predchádzajúci kód
a informácia o chybe, aby mohol model vytvoriť opravenú verziu.

\subsection{Izolované vykonávanie Python kódu}

Vygenerovaný Python kód sa nevykonáva priamo v hlavnom procese programu.
Na jeho vykonanie sa používa Docker kontajner. Tento prístup umožňuje oddeliť
vykonávanie kódu generovaného jazykovým modelom od hlavnej aplikácie.

Kontajner sa spúšťa s viacerými obmedzeniami. Sieťová komunikácia je
zakázaná pomocou parametra \texttt{--network none}. Kontajner má zároveň
obmedzenú operačnú pamäť na 512 MB a využitie procesora na jeden CPU.
Počet procesov je obmedzený pomocou parametra \texttt{--pids-limit}.

Súborový systém kontajnera je nastavený ako iba na čítanie. Pre dočasné
operácie je povolený samostatný adresár \texttt{/tmp}. Zároveň sú odstránené
Linuxové capabilities pomocou \texttt{--cap-drop ALL} a je použitá možnosť
\texttt{no-new-privileges}.

Do kontajnera sa pripájajú iba potrebné súbory. Ide o príslušnú tabuľku
\texttt{table.parquet}, súbor so vygenerovaným kódom
\texttt{generated\_code.py} a súbor \texttt{sandbox\_runner.py}. Všetky
tieto súbory sú pripojené v režime iba na čítanie.

Samotné vykonanie je časovo obmedzené na 20 sekúnd. Ak program prekročí
tento limit alebo Docker vráti chybu, vykonanie sa považuje za neúspešné.

\subsection{Spracovanie výsledku v sandboxe}

Po spustení kontajnera načíta súbor \texttt{sandbox\_runner.py} príslušnú
tabuľku do objektu \texttt{df}. Vygenerovaný Python kód sa následne vykoná
v priestore, v ktorom je objekt \texttt{df} dostupný.

Po vykonaní kódu sa kontroluje, či vygenerovaný program vytvoril premennú
\texttt{result}. Ak táto premenná neexistuje, vykonanie sa považuje za
neúspešné.

Výsledok sa následne normalizuje pomocou funkcie
\texttt{normalize\_result()}. Táto funkcia konvertuje hodnoty z knižníc
Pandas a NumPy na štandardné Python dátové typy, ktoré je možné serializovať
do formátu JSON. Podporované sú napríklad skalárne hodnoty, zoznamy,
slovníky a hodnoty získané z objektov \texttt{Series} alebo
\texttt{ndarray}.

Takto pripravený výsledok je následne odoslaný späť do hlavného programu
vo forme JSON. Hlavný program ho dekóduje a používa ako predikovanú odpoveď
na danú otázku.

\subsection{Opakovanie po chybe vykonávania}

Ak sa vygenerovaný kód nepodarí vykonať v Docker kontajneri, systém sa
pokúsi vytvoriť nový kód. Maximálne sú teda použité dva pokusy v rámci
spracovania jednej otázky.

Pri druhom pokuse sa do promptu pridá informácia o predchádzajúcom
neúspešnom vykonaní a predchádzajúci vygenerovaný kód. Jazykový model je
požiadaný o vytvorenie opravenej verzie.

Aktuálna implementácia používa pri chybe Dockeru všeobecnú informáciu
o neúspešnom vykonaní. Konkrétny traceback z Docker kontajnera sa síce
vypíše do konzoly, ale ešte nie je štruktúrovane odovzdaný späť jazykovému
modelu.

\subsection{Vyhodnotenie odpovede}

Po úspešnom vykonaní Python kódu sa získaná hodnota porovná s referenčnou
odpoveďou. Na tento účel sa používa trieda \texttt{Evaluator} z balíka
\texttt{databench\_eval}.

Funkcia \texttt{evaluate\_answer()} využíva porovnávaciu funkciu
\texttt{default\_compare()}, ktorá zohľadňuje aj typ odpovede. Výsledkom
porovnania je logická hodnota určujúca, či sa predikovaná odpoveď považuje
za správnu.

V režime \texttt{full} sa ako referenčná odpoveď používa hodnota
\texttt{answer}. V režime \texttt{lite} sa používa hodnota
\texttt{sample\_answer}. V prípade, že referenčná odpoveď chýba, otázka
sa označí stavom \texttt{unscored} a nepovažuje sa za úspešne vyhodnotenú.

\subsection{Ukladanie výsledkov}

Výsledky jednotlivých otázok sa ukladajú do súboru
\texttt{results.jsonl}. Každý riadok obsahuje jeden JSON objekt.

Medzi ukladané informácie patria napríklad index otázky, identifikátor
tabuľky, režim spracovania, samotná otázka, typ odpovede, očakávaná
odpoveď, predikovaná odpoveď, informácia o správnosti, stav spracovania,
chybová správa, vygenerovaný kód a počet pokusov.

Takéto uloženie umožňuje neskoršie analyzovať jednotlivé výsledky bez
potreby opätovného spúšťania jazykového modelu.

\subsection{Režimy spustenia programu}

Program \texttt{main.py} podporuje dva režimy spracovania otázok:
\texttt{single} a \texttt{batch}.

Režim \texttt{single} umožňuje vybrať jednu otázku podľa jej indexu.
Tento režim je vhodný najmä na manuálne overovanie správania systému,
promptu a generovaného kódu.

Režim \texttt{batch} umožňuje spracovať súvislý rozsah otázok. Používateľ
zadá počiatočný index a počet otázok. Program následne postupne spracuje
jednotlivé otázky a priebežne ukladá ich výsledky do súboru
\texttt{results.jsonl}.

Oba režimy podporujú režim dát \texttt{full} a \texttt{lite}. Týmto
spôsobom je možné experimentálne porovnať spracovanie úplných tabuliek
a menších vzoriek.

\subsection{Testovanie implementácie}

Súčasťou projektu sú aj testovacie skripty, ktoré overujú jednotlivé časti
implementácie. Testovanie zahŕňa napríklad kontrolu načítania lokálnych
tabuliek, správanie evaluátora, spracovanie chýb Docker kontajnera a
správanie opakovania po chybe.

Pre kontrolu syntaktickej správnosti Python súborov je možné použiť
modul \texttt{py\_compile}. Samostatný test
\texttt{test\_retry.py} je implementovaný ako štandardný
\texttt{unittest} test pre kontrolu opakovania po neúspešnom vykonaní.
Ostatné testovacie súbory sú v aktuálnej verzii projektu pripravené
predovšetkým ako manuálne validačné skripty.

Implementované riešenie realizuje základný baseline založený na generovaní
Python kódu jazykovým modelom. Celý proces pozostáva z načítania otázky
a tabuľky, vytvorenia schémy a promptu, generovania Python kódu, kontroly
jeho syntaxe, izolovaného vykonania v Docker kontajneri, normalizácie
výsledku a jeho porovnania s referenčnou odpoveďou.

Navrhnutá architektúra zároveň umožňuje zaznamenávať výsledky a chyby
jednotlivých otázok. To poskytuje základ pre experimentálne vyhodnotenie
baseline a pre ďalšie rozšírenia projektu v nasledujúcej etape.

\section{Metodika hodnotenia a použité metriky}
\label{sec:metodika-hodnotenia}

\subsection{Cieľ hodnotenia}

Cieľom hodnotenia je overiť, do akej miery dokáže navrhovaný systém správne odpovedať na otázky nad tabuľkovými dátami v rámci benchmarku DataBench. Hodnotenie sa zameriava na správnosť výsledných odpovedí a na schopnosť systému úspešne dokončiť proces generovania a vykonávania Python kódu.

Keďže navrhovaný baseline využíva jazykový model na generovanie kódu, úspešnosť systému závisí od viacerých faktorov. Model musí správne interpretovať otázku, identifikovať relevantné stĺpce tabuľky, zvoliť vhodný spôsob spracovania údajov a vytvoriť odpoveď v požadovanom formáte. Chyba v ktorejkoľvek z týchto fáz môže viesť k nesprávnemu výsledku alebo k neúspešnému dokončeniu úlohy.

Na vyhodnotenie odpovedí sa využíva nástroj
\texttt{databench\_eval}, ktorý porovnáva predikcie s referenčnými odpoveďami podľa pravidiel benchmarku. Výsledky spracovania sa zároveň ukladajú do súboru \texttt{results.jsonl}, ktorý obsahuje záznamy o jednotlivých otázkach a ich vyhodnotení.

\subsection{Metriky hodnotenia}

Implementácia baseline používa dve základné metriky: \textit{Accuracy} a \textit{Execution Error Rate}. Pri interpretácii týchto metrík je potrebné zohľadniť spôsob, akým aktuálna implementácia rozlišuje úspešne spracované otázky od ostatných prípadov.

\subsubsection{Accuracy}

Metrika \textit{Accuracy} vyjadruje podiel správnych odpovedí spomedzi otázok, ktoré systém úspešne spracoval. V aktuálnej implementácii sa počet správnych odpovedí delí počtom úspešne spracovaných otázok.

Metrika sa vypočíta podľa vzorca:

\begin{equation}
\mathrm{Accuracy} =
\frac{N_{\mathrm{correct}}}
{N_{\mathrm{successful}}}
\end{equation}

kde $N_{\mathrm{correct}}$ predstavuje počet správnych odpovedí a
$N_{\mathrm{successful}}$ počet úspešne spracovaných otázok.

Správnosť odpovedí sa určuje pomocou hodnotiaceho nástroja
\texttt{databench\_eval}. Pri hodnotení je potrebné rešpektovať očakávaný typ odpovede, pretože benchmark obsahuje rôzne formáty výstupov, napríklad logické hodnoty, kategórie, číselné hodnoty a zoznamy hodnôt.

Keďže sa pri výpočte tejto metriky nezohľadňujú neúspešne spracované otázky v menovateli, výslednú hodnotu je potrebné interpretovať spolu s počtom úspešne spracovaných otázok a počtom všetkých spracovaných otázok.

\subsubsection{Execution Error Rate}

Metrika \textit{Execution Error Rate} v aktuálnej implementácii vyjadruje podiel otázok, ktoré neboli označené ako úspešne spracované, z celkového počtu spracovaných otázok.

V implementácii sa vypočíta podľa vzorca:

\begin{equation}
\mathrm{ExecutionErrorRate} =
\frac{N_{\mathrm{total}} - N_{\mathrm{successful}}}
{N_{\mathrm{total}}}
\end{equation}

kde $N_{\mathrm{total}}$ označuje celkový počet spracovaných otázok a
$N_{\mathrm{successful}}$ počet otázok označených ako úspešne spracované.

Táto metrika poskytuje prehľad o podiele prípadov, pri ktorých sa nepodarilo úspešne dokončiť spracovanie otázky. Vzhľadom na aktuálnu implementáciu však nejde výlučne o mieru chýb pri vykonávaní Python kódu. Medzi neúspešné prípady môžu patriť aj iné stavy spracovania, ktoré nie sú označené ako úspešné.

Preto sa táto hodnota interpretuje ako ukazovateľ neúspešného spracovania otázok, nie ako presný počet chýb jednotlivých vykonaní kódu. Na samostatné meranie chýb vykonávania by bolo potrebné evidovať počet pokusov o vykonanie kódu a počet neúspešných vykonaní nezávisle od ostatných stavov.

\subsubsection{Počet úspešne spracovaných otázok}

Okrem uvedených metrík sa sleduje aj počet otázok, ktoré systém úspešne spracoval. Tento údaj umožňuje posúdiť, na akom počte otázok je založená hodnota metriky \textit{Accuracy}.

Samotné úspešné spracovanie otázky ešte nezaručuje správnosť odpovede. Preto je potrebné tento údaj interpretovať spolu s metrikou \textit{Accuracy} a s celkovým počtom spracovaných otázok.

\subsection{Postup hodnotenia}

Hodnotenie pozostáva z niekoľkých krokov. Najskôr systém načíta otázku a identifikátor príslušnej tabuľky. Následne načíta tabuľkové dáta a pomocou jazykového modelu vytvorí Python kód určený na spracovanie týchto dát.

Vygenerovaný kód sa následne vykoná v kontrolovanom prostredí. Výsledok sa spracuje a prevedie do očakávaného formátu odpovede. Počas tohto procesu sa zaznamenávajú aj prípady, pri ktorých sa nepodarí úspešne dokončiť spracovanie otázky.

Výsledky jednotlivých otázok sa ukladajú do súboru
\texttt{results.jsonl}. Tento súbor slúži ako záznam o priebehu spracovania a obsahuje údaje potrebné na kontrolu výsledkov a analýzu neúspešných prípadov.

Na vyhodnotenie správnosti predikcií sa používa nástroj
\texttt{databench\_eval}, ktorý porovnáva vytvorené odpovede s referenčnými odpoveďami podľa pravidiel benchmarku.

Pri analýze výsledkov je potrebné rozlišovať medzi úspešným vykonaním kódu a správnosťou odpovede. Aj v prípade, keď sa kód vykoná bez chyby, môže vytvorená odpoveď obsahovať nesprávnu hodnotu. Naopak, ak sa spracovanie otázky nepodarí dokončiť, prípad sa musí zohľadniť pri interpretácii celkovej úspešnosti systému.

Aktuálna implementácia zaznamenáva výsledky vo formáte
\texttt{results.jsonl}. Samostatné súbory
\texttt{predictions.txt} a \texttt{predictions\_lite.txt}, ktoré by obsahovali iba predikcie v poradí otázok požadovanom benchmarkom, nie sú v súčasnej verzii implementované.

\subsection{Porovnanie DataBench a DataBench Lite}

Pri vyhodnocovaní je potrebné rozlišovať medzi variantmi DataBench a DataBench Lite. Plná verzia pracuje s úplnými tabuľkami, zatiaľ čo variant Lite využíva vzorky tabuliek.

Rozsah dostupných údajov môže ovplyvniť schopnosť systému správne odpovedať na otázky. Pri práci so vzorkami nemusia byť dostupné všetky riadky potrebné na získanie správnej odpovede. Výsledky oboch variantov preto treba interpretovať s ohľadom na použitý rozsah dát.

Pre korektné porovnanie je potrebné vyhodnotiť oba varianty samostatne a zaznamenať použitú konfiguráciu experimentu. Zároveň je potrebné zabezpečiť, aby boli výsledky priradené k správnemu variantu datasetu a aby sa pri ich interpretácii zohľadnili rozdiely v použitých dátach.

\subsection{Reprodukovateľnosť experimentov}

Na zabezpečenie reprodukovateľnosti je potrebné zaznamenávať parametre použité pri jednotlivých experimentoch. Patria medzi ne najmä identifikácia jazykového modelu, konfigurácia generovania, verzia promptu, použitá verzia hodnotiaceho nástroja a počet spracovaných otázok. Konfigurácia použitá v aktuálnej implementácii baseline je uvedená v tabuľke~\ref{tab:repro}.

\% V preambule:
% \usepackage{booktabs}
% \usepackage[table]{xcolor}
% \usepackage{array}
% \definecolor{groupgray}{gray}{0.92}

\begin{table}[h]
\centering
\small
\caption{Konfigurácia experimentov baseline}
\label{tab:repro}
\renewcommand{\arraystretch}{1.25}
\begin{tabular}{@{}>{\raggedright\arraybackslash}p{4.6cm} >{\raggedright\arraybackslash}p{8.6cm}@{}}
\toprule
\textbf{Oblasť} & \textbf{Použité nastavenie} \\
\midrule

\multicolumn{2}{@{}l}{\textbf{Generovanie kódu}} \\
Model & \texttt{gemini-3.5-flash-lite} cez Google Gemini API (knižnica \texttt{google-genai}, verzia: \textit{doplniť}) \\
Parametre generovania & predvolené hodnoty API; \texttt{temperature}, \texttt{top\_p}, \texttt{top\_k} ani \texttt{seed} nie sú nastavené, preto generovanie nie je deterministické \\
Prompt & verzia v1 (\texttt{build\_prompt()}), bez systémového promptu; obsahuje otázku, schému tabuľky, očakávaný typ odpovede a pravidlá pre výstup \\
Počet pokusov & najviac 2 generovania na otázku (prvotné a jedno opravné); pri limite požiadaviek API až 5 opakovaní s čakaním 5--60~s \\

\addlinespace
\multicolumn{2}{@{}l}{\textbf{Vykonávanie kódu}} \\
Prostredie & Docker, obraz \texttt{python:3.13-slim} s knižnicami \texttt{pandas} a \texttt{pyarrow} (verzie nie sú pripnuté) \\
Obmedzenia & bez siete, 512~MB RAM, 1~CPU, časový limit 20~s \\

\addlinespace
\multicolumn{2}{@{}l}{\textbf{Dáta a vyhodnotenie}} \\
Dataset & \texttt{cardiffnlp/databench}, konfigurácia \texttt{qa}, rozdelenie \texttt{train}; počet záznamov a revízia: \textit{doplniť} \\
Režimy & \texttt{full}: \texttt{all.parquet}, referenčná odpoveď \texttt{answer}; \texttt{lite}: \texttt{sample.parquet}, referenčná odpoveď \texttt{sample\_answer} \\
Hodnotenie & \texttt{databench\_eval} (verzia alebo commit: \textit{doplniť}) \\
Ostatné knižnice & \texttt{datasets}, \texttt{pandas}, \texttt{pyarrow}, \texttt{numpy} (verzie: \textit{doplniť}) \\

\bottomrule
\end{tabular}
\end{table}

Keďže teplota ani náhodné semienko (\texttt{seed}) nie sú explicitne nastavené, generovanie kódu nie je deterministické. Opakované spustenie na rovnakých otázkach preto môže vytvoriť odlišný kód a mierne odlišné výsledky. Z tohto dôvodu sa pri každom experimente ukladá vygenerovaný kód aj predikcia, aby bolo možné výsledky spätne skontrolovať.

Súbor \texttt{results.jsonl} poskytuje záznamy o spracovaní jednotlivých otázok a umožňuje spätne analyzovať výsledky. Každý záznam obsahuje režim dát (\texttt{mode}), vygenerovaný kód a počet pokusov (\texttt{attempts}). Na úplnú reprodukovateľnosť je však potrebné zaznamenať aj konfiguráciu experimentu a verzie použitých knižníc a nástrojov, ktoré sú uvedené vyššie.

Pri interpretácii výsledkov sa musí uvádzať aj použitá množina dát. Výsledky získané na tréningových alebo vývojových dátach sa nesmú zamieňať s výsledkami záverečného hodnotenia na testovacích dátach. Referenčné odpovede z rozdelenia \texttt{qa}/\texttt{train}, ktoré sa používajú v tejto etape, slúžia iba na vývojovú kontrolu baseline.

Referenčné testovacie odpovede sa nemajú používať na priebežné prispôsobovanie promptov ani na ladenie systému, pretože takýto postup by mohol ovplyvniť objektivitu záverečného hodnotenia.

\subsection{Prezentácia výsledkov}

Výsledky hodnotenia baseline sa sumarizujú pomocou počtu spracovaných otázok, počtu úspešne spracovaných otázok a hodnôt metrík \textit{Accuracy} a \textit{Execution Error Rate}. Tabuľka~\ref{tab:baseline-metriky} poskytuje priestor na porovnanie výsledkov jednotlivých variantov benchmarku.

\begin{table}[htbp]
    \centering
    \caption{Vyhodnotenie baseline}
    \label{tab:baseline-metriky}
    \begin{tabular}{|l|c|c|}
        \hline
        \textbf{Metrika} &
        \textbf{DataBench} &
        \textbf{DataBench Lite} \\
        \hline
        Počet spracovaných otázok & -- & -- \\
        \hline
        Počet úspešne spracovaných otázok & -- & -- \\
        \hline
        Počet správnych odpovedí & -- & -- \\
        \hline
        Accuracy & -- & -- \\
        \hline
        Execution Error Rate & -- & -- \\
        \hline
    \end{tabular}
\end{table}

Hodnoty v tabuľke sa doplnia po vykonaní experimentov. Pri prezentácii výsledkov je potrebné uviesť, že hodnota \textit{Accuracy} sa v aktuálnej implementácii počíta iba z úspešne spracovaných otázok, zatiaľ čo \textit{Execution Error Rate} predstavuje podiel otázok, ktoré neboli označené ako úspešne spracované.

Takto prezentované výsledky umožňujú posúdiť správnosť odpovedí aj schopnosť systému dokončiť spracovanie otázok. Zároveň poskytujú základ na identifikáciu obmedzení aktuálneho baseline a na plánovanie jeho ďalších zlepšení.


\section{Obmedzenia riešenia a plán ďalšej etapy}

Implementovaný baseline slúži ako základ pre experimenty v druhej etape projektu. V tejto kapitole sú zhrnuté známe obmedzenia aktuálnej implementácie, ich možný dopad na výsledky a plánované úpravy. Ide o plán práce a analýzu rizík, nie o už dosiahnuté výsledky. Obmedzenia vychádzajú z analýzy zdrojového kódu a z priebehu doterajšieho testovania.

\subsection{Obmedzenia aktuálnej implementácie}

\subsubsection*{Generovanie kódu a prompt}

\begin{itemize}
    \item \textbf{Obmedzená informácia o dátach v prompte.} Model dostáva iba názvy stĺpcov, ich dátové typy a očakávaný typ odpovede. Nevidí ukážkové riadky ani konkrétne hodnoty v kategoriálnych stĺpcoch. Pri otázkach, kde záleží na presnom zápise hodnoty (napríklad názov krajiny alebo kategórie), môže model použiť nesprávny reťazec a filter nevráti žiadne riadky.
    \item \textbf{Jediná verzia promptu.} Prompt obsahuje všeobecné pravidlá, ale neobsahuje ukážky riešení (few-shot) a ani sa nemení podľa typu otázky okrem krátkych inštrukcií pre jednotlivé typy odpovede. Vplyv jednotlivých pravidiel na presnosť nebol zatiaľ meraný.
    \item \textbf{Jazyk otázok.} Dataset obsahuje otázky v angličtine aj španielčine, zatiaľ však nebolo overené, ako sa prístup správa pri španielskych otázkach.
    \item \textbf{Nedeterminizmus generovania.} Teplota ani \texttt{seed} nie sú nastavené. Opakované spustenie na rovnakých otázkach môže priniesť odlišný kód a mierne odlišné výsledky, čo komplikuje porovnanie experimentov.
\end{itemize}

\subsubsection*{Opakovanie po chybe}

\begin{itemize}
    \item \textbf{Spätná väzba pri chybách.} Pri zlyhaní vykonania v Dockeri funkcia \texttt{run\_code\_in\_docker()} vráti iba hodnotu \texttt{None} a model dostane všeobecnú informáciu o neúspechu. Konkrétny traceback sa vypíše do konzoly, ale nie je štruktúrovane odovzdaný späť, takže oprava v druhom pokuse je menej cielená.
    \item \textbf{Nerozlíšené príčiny zlyhania.} Hodnota \texttt{None} spája rôzne situácie: chybu v kóde, prekročenie časového limitu, nedostatok pamäte aj nedostupnosť Dockeru. Z uložených záznamov preto nie je možné spoľahlivo určiť, ktorá príčina zlyhania prevláda.
    \item \textbf{Iba jedna oprava.} V rámci jednej otázky je povolený jeden opravný pokus. Zložitejšie otázky môžu vyžadovať viac iterácií, no zatiaľ nebolo skúmané, či by ďalšie pokusy zlepšili výsledok.
\end{itemize}

\subsubsection*{Vykonávanie kódu}

\begin{itemize}
    \item \textbf{Zdieľaný súbor s kódom.} Vygenerovaný kód sa ukladá do jedného súboru \texttt{generated\_code.py}. Súbežné spustenia by sa navzájom prepisovali, preto spracovanie otázok prebieha sekvenčne a nie je ho možné bezpečne paralelizovať bez samostatných dočasných súborov.
    \item \textbf{Výkon.} Pre každú otázku sa spúšťa nový Docker kontajner. Pri väčších dávkach otázok tak vzniká réžia spúšťania a celkový čas spracovania je dlhší.
    \item \textbf{Zdrojové limity.} Limit 512~MB pamäte a 20~s času je vhodný na ochranu hostiteľského systému, ale pri veľkých tabuľkách v režime \texttt{full} môže viesť k zlyhaniu aj pri správnom kóde. Rozsah tohto problému ešte nebol vyhodnotený.
    \item \textbf{Bezpečnosť.} Docker kontajner s obmedzeniami (bez siete, odobraté oprávnenia, súborový systém iba na čítanie) znižuje riziko, ale nie je formálnou zárukou bezpečnosti pri vykonávaní kódu generovaného jazykovým modelom.
\end{itemize}

\subsubsection*{Vyhodnotenie a metriky}

\begin{itemize}
    \item \textbf{Accuracy iba z úspešných otázok.} Metrika sa počíta ako podiel správnych odpovedí z úspešne spracovaných otázok. Otázky, ktoré sa nepodarilo dokončiť, sa do menovateľa nezahŕňajú, a hodnota môže byť preto optimistickejšia než presnosť počítaná zo všetkých otázok.
    \item \textbf{Hrubá kategorizácia stavov.} Execution Error Rate zahŕňa všetky neúspešné stavy vrátane \texttt{unscored}, ktoré nie sú chybami vykonania. Neexistuje samostatné rozlíšenie chýb generovania, syntaxe, vykonania a nesprávnej odpovede.
    \item \textbf{Zmiešané režimy.} Metriky sa počítajú zo všetkých záznamov v súbore \texttt{results.jsonl}. Ak súbor obsahuje výsledky oboch režimov, nie sú vyhodnotené samostatne.
    \item \textbf{Závislosť od internej funkcie.} Porovnanie odpovedí využíva metódu \texttt{default\_compare()} triedy \texttt{Evaluator} z balíka \texttt{databench\_eval}, pričom objekt sa vytvára bez štandardnej inicializácie. Zmena v budúcej verzii balíka by mohla toto správanie ovplyvniť, preto je dôležité pripnúť jeho verziu.
\end{itemize}

\subsubsection*{Dáta, výstup a prostredie}

\begin{itemize}
    \item \textbf{Vývojové dáta.} Experimenty používajú konfiguráciu \texttt{qa} a rozdelenie \texttt{train}, teda slúžia na vývojovú kontrolu. Nejde o oficiálne vyhodnotenie na testovacích dátach súťaže.
    \item \textbf{Chýbajúci súťažný výstup.} Súbory \texttt{predictions.txt} a \texttt{predictions\_lite.txt} v poradí otázok požadovanom benchmarkom nie sú implementované.
    \item \textbf{Závislosť od externého API.} Dostupnosť služby Gemini ovplyvňuje priebeh experimentov. Pri testovaní sa vyskytla dočasná chyba \texttt{503 UNAVAILABLE}, ktorá nesúvisí s logikou riešenia. Opakovanie požiadavky pri takýchto chybách je potrebné ošetriť a otestovať.
    \item \textbf{Reprodukovateľnosť prostredia.} Verzie knižníc nie sú pripnuté v súbore \texttt{requirements.txt} a v Docker obraze nie sú pripnuté verzie \texttt{pandas} a \texttt{pyarrow}.
\end{itemize}

\subsection{Zhrnutie obmedzení a plánovaných úprav}

Tabuľka~\ref{tab:limits} prehľadne spája najdôležitejšie obmedzenia s ich možným dopadom a plánovanou úpravou.

\begin{table}[h]
\centering
\small
\caption{Obmedzenia baseline a plánované úpravy}
\label{tab:limits}
\renewcommand{\arraystretch}{1.25}
\begin{tabular}{@{}>{\raggedright\arraybackslash}p{4.2cm} >{\raggedright\arraybackslash}p{4.6cm} >{\raggedright\arraybackslash}p{4.6cm}@{}}
\toprule
\textbf{Obmedzenie} & \textbf{Možný dopad} & \textbf{Plánovaná úprava} \\
\midrule
Všeobecná informácia o chybe v druhom pokuse & menej cielená oprava kódu & odovzdať modelu štruktúrovaný traceback \\
Prompt bez ukážok hodnôt & nesprávne filtre pri kategoriálnych stĺpcoch & experiment s pridaním ukážkových riadkov alebo hodnôt \\
Accuracy iba z úspešných otázok & optimistický odhad presnosti & počítať Accuracy aj zo všetkých otázok \\
Nerozlíšené stavy a režimy & neprehľadná analýza chýb & samostatné stavy a metriky pre \texttt{full} a \texttt{lite} \\
Chýbajúce \texttt{predictions.txt} & nemožnosť odovzdať súťažný výstup & implementovať export oboch súborov \\
Nedeterministické generovanie & ťažšie porovnanie experimentov & nastaviť teplotu a \texttt{seed} \\
Zdieľaný \texttt{generated\_code.py} & nemožnosť paralelizácie & jedinečné dočasné súbory pre každé spustenie \\
Nepripnuté verzie knižníc & ťažšia reprodukovateľnosť & \texttt{requirements.txt}, pripnutý Docker obraz \\
\bottomrule
\end{tabular}
\end{table}

\subsection{Plán ďalšej etapy}

V druhej etape projektu plánujeme postupovať v nasledujúcich krokoch. Poradie zohľadňuje, že najprv je potrebné zabezpečiť spoľahlivé meranie a až potom experimentovať so zlepšeniami.

\paragraph{1. Úprava merania.} Pred porovnávaním variantov riešenia upravíme metriky tak, aby rozlišovali chyby generovania, syntaxe a vykonania, nesprávne odpovede a stav \texttt{unscored}, a aby sa režimy \texttt{full} a \texttt{lite} vyhodnocovali samostatne. Okrem súčasnej Accuracy bude uvedená aj hodnota zo všetkých otázok. Tým sa zabezpečí, že neskoršie zlepšenia bude možné posudzovať na spoľahlivom základe.

\paragraph{2. Export predikcií.} Implementujeme export súborov \texttt{predictions.txt} a \texttt{predictions\_lite.txt} v poradí otázok podľa testovacieho súboru, aby bolo riešenie pripravené na oficiálne vyhodnotenie.

\paragraph{3. Analýza chýb.} Nesprávne odpovede rozdelíme podľa typu odpovede (\texttt{boolean}, \texttt{number}, \texttt{category}, \texttt{list[category]}, \texttt{list[number]}) a podľa príčiny: nesprávny výber stĺpca, nesprávna interpretácia otázky, nesprávny formát výstupu alebo chyba vo výpočte. Z toho určíme, ktorým úpravám dať prednosť.

\paragraph{4. Experimenty s promptom a opakovaním.} Porovnáme viacero verzií promptu, napríklad s ukážkovými riadkami a hodnotami kategoriálnych stĺpcov alebo s ukážkami riešení, a overíme vplyv štruktúrovaného tracebacku a väčšieho počtu opravných pokusov. Každá zmena sa bude testovať na rovnakej množine otázok a pri rovnakej konfigurácii.

\paragraph{5. Porovnanie s priamym generovaním odpovede.} Ako doplnkový experiment porovnáme generovanie Python kódu s priamym generovaním odpovede jazykovým modelom bez vykonania kódu.

\paragraph{6. Vyhodnotenie na súťažných rozdeleniach.} Riešenie overíme aj na rozdeleniach \texttt{dev} a \texttt{test} konfigurácie \texttt{semeval}. Referenčné testovacie odpovede sa nebudú používať na priebežné ladenie promptov.

\paragraph{7. Technická stabilita a reprodukovateľnosť.} Zavedieme jedinečné dočasné súbory pre každé spustenie, pripneme verzie knižníc a Docker obrazu, nastavíme teplotu a \texttt{seed} a ošetríme opakovanie požiadaviek pri dočasných chybách API.

Poradie a rozsah týchto úprav sa môže zmeniť podľa výsledkov analýzy chýb. Úpravy, ktoré sa nepodarí vykonať v dostupnom čase, budú uvedené ako návrhy do budúcnosti.

\section{Použitie AI nástrojov}

V projekte sa vyskytujú dva odlišné spôsoby použitia umelej inteligencie. Prvým je jazykový model Gemini, ktorý je súčasťou samotného riešenia. Druhým sú AI asistenti, ktoré členovia tímu využívali ako pomocný nástroj pri vývoji a tvorbe dokumentácie. Obe oblasti sú nižšie opísané samostatne, aby bolo zrejmé, ktoré výsledky pochádzajú z behu systému a ktoré texty vznikli s podporou AI asistentov.

\subsection{Gemini ako súčasť systému}

Jazykový model \texttt{gemini-3.5-flash-lite} (Google Gemini API) nie je pomocným nástrojom vývoja. Je to funkčná súčasť implementovaného baseline: pre každú otázku generuje Python kód, ktorý sa následne vykoná v izolovanom prostredí a jeho výsledok sa vyhodnotí. Konfigurácia modelu je uvedená v tabuľke~\ref{tab:repro} a postup spracovania v kapitole~4.

Zloženie promptu určuje funkcia \texttt{build\_prompt()} v súbore \texttt{prompt\_builder.py}. Prompt obsahuje otázku, schému tabuľky, očakávaný typ odpovede a pravidlá pre výstup, napríklad:

\begin{quote}
\small\itshape
You are a Python data analysis assistant. A pandas DataFrame named df is already loaded. [\ldots] Store the final answer in a variable named result. [\ldots] Return only executable Python code, without explanations or markdown.
\end{quote}

Vygenerovaný kód, predikcia a počet pokusov sa pre každú otázku ukladajú do súboru \texttt{results.jsonl}. Činnosť Gemini je teda priamo súčasťou meraných výsledkov a nie je skrytou pomocou pri vývoji.

\subsection{AI nástroje použité pri vývoji a dokumentácii}

Pri vývoji a písaní dokumentácie boli AI nástroje využívané iba ako pomocný nástroj. Prehľad použitých nástrojov a účelu ich použitia je v tabuľke~\ref{tab:ai-tools}.

\begin{table}[h]
\centering
\small
\caption{Prehľad AI nástrojov použitých pri vývoji a dokumentácii}
\label{tab:ai-tools}
\renewcommand{\arraystretch}{1.25}
\begin{tabular}{@{}>{\raggedright\arraybackslash}p{2.6cm} >{\raggedright\arraybackslash}p{5.2cm} >{\raggedright\arraybackslash}p{5.4cm}@{}}
\toprule
\textbf{Nástroj} & \textbf{Účel použitia} & \textbf{Typ výstupu} \\
\midrule
ChatGPT & vysvetlenie technických konceptov, pomoc pri pochopení vlastného kódu, konzultácia možných spôsobov implementácie a štruktúry riešenia, jazyková kontrola textu, pomoc pri organizácii dokumentácie & vysvetlenia, návrhy štruktúry, jazykové úpravy textu \\
Claude (Anthropic) & kontrola súladu dokumentácie s kódom projektu, identifikácia nedostatkov dokumentácie, pomoc pri analýze chybového výpisu, formulácia a formátovanie vybraných častí textu v \LaTeX{}u (konfigurácia experimentov, táto kapitola) & zoznam pripomienok k dokumentácii, diagnóza chyby, návrhy textu \\
\bottomrule
\end{tabular}
\end{table}

Využitie AI asistentov zahŕňalo najmä:

\begin{itemize}
    \item vysvetlenie vybraných technických konceptov súvisiacich s riešením,
    \item pomoc pri analýze a pochopení častí vlastného kódu,
    \item konzultáciu možných spôsobov implementácie a štruktúry riešenia,
    \item kontrolu súladu dokumentácie s implementáciou,
    \item pomoc pri hľadaní príčin chýb počas programovania,
    \item kontrolu a zlepšenie formulácie textu v projektovej dokumentácii,
    \item pomoc pri organizácii a štruktúrovaní jednotlivých častí dokumentácie.
\end{itemize}

\subsection{Ukážky použitia asistenta Claude}

Nižšie sú uvedené príklady použitia asistenta Claude, pri ktorých máme záznam rozhovoru. Prompty boli pôvodne napísané v ukrajinčine, v texte sú preložené do slovenčiny. Výstupy sú zhrnuté, nie citované celé.

\paragraph{Príklad 1 -- porovnanie dokumentácie s riešením (Claude).}
\emph{Prompt:} Po nahraní PDF dokumentácie: „Máme takúto dokumentáciu. Je jej gramatika správna?“\\
\emph{Výstup:} Zoznam nedostatkov, napríklad prázdna tabuľka výsledkov, neplatný odkaz \texttt{[?]} v kapitole~3.1, chýbajúci presný identifikátor modelu a príliš stručná kapitola o AI nástrojoch.\\
\emph{Použitie tímom:} Pripomienky sme posúdili a zapracovali do tejto verzie dokumentu.

\paragraph{Príklad 2 -- formulácia dokumentácie (Claude).}
\emph{Prompt:} „Pridajte presný názov modelu, teplotu, verziu promptu a verzie knižníc.“\\
\emph{Výstup:} Návrh textu kapitoly o reprodukovateľnosti a tabuľky konfigurácie. Hodnoty asistent prevzal zo zdrojového kódu a uviedol, že teplota nie je nastavená, a preto sa nemá uvádzať ako nulová.\\
\emph{Použitie tímom:} Hodnoty sme porovnali s kódom a chýbajúce údaje (verzie knižníc) doplnili členovia tímu.

\paragraph{Príklad 3 -- analýza chyby pri programovaní (GPT).}
\emph{Prompt:} Vložený výpis chyby z terminálu (\texttt{503 UNAVAILABLE}, \emph{This model is currently experiencing high demand}) so slovami „znova vznikla chyba, pomoc s testovaním“.\\
\emph{Výstup:} Asistent určil, že chyba nevzniká v logike projektu, ale z dôvodu dočasného preťaženia serverov modelu. Zároveň upozornil na dva problémy v kóde. Funkcia \texttt{generate\_code()} opakovala požiadavku iba pri chybe 429, takže pri chybe 503 by spracovanie otázky skončilo hneď. Testovacie volanie API bolo v súbore \texttt{llm\_client.py} mimo bloku \texttt{if \_\_name\_\_ == "\_\_main\_\_"}, a preto by sa vykonávalo pri každom importe modulu. Navrhol rozšíriť opakovanie o chyby 503 a \texttt{UNAVAILABLE} a pripravil unit test s nahradeným klientom, ktorý nevyžaduje sieť.\\
\emph{Použitie tímom:} Návrh sme posúdili ako podklad na úpravu kódu. Diagnózu je možné overiť z priloženého výpisu chyby.

\paragraph{Príklad 4 -- Pomoc pri testovaní baselinu (GPT).}
\emph{Prompt:} Popis fungovania projektu a žiadosť o vyznačenie rizikových miest, ktoré je potrebné otestovať\\
\emph{Výstup:} Popis rizikových miest a podrobných postupov ich testovania \\
\emph{Použitie tímom:} Identifikácia kritických miest a testovanie sa vykonali pomocou Chat GPT.

\paragraph{Príklad 5 -- Hodnotenie výsledkov testovania (GPT).}
\emph{Prompt:} Popis získaných výsledkov testovania. Prosím, ohodnoťte ich a navrhnite, ako by sa dali zlepšiť.\\
\emph{Výstup:} Hodnotenie výsledkov ako dobré a poskytnutie niekoľkých návrhov na vylepšenie kódu s cieľom dosiahnuť lepšie výsledky \\
\emph{Použitie tímom:} Presvedčenie o dobrých výsledkoch a pochopenie smeru zlepšovania kódu v druhej fáze.

AI nástroje neboli použité na automatické vytvorenie kompletného riešenia ani na nahradenie vlastnej implementácie. Návrh riešenia, písanie a úprava zdrojového kódu, spúšťanie experimentov a interpretácia výsledkov zostali zodpovednosťou členov tímu. Texty a návrhy získané od AI asistentov boli pred zapracovaním porovnané so skutočnou implementáciou a upravené tak, aby zodpovedali stavu projektu. Merané výsledky pochádzajú z behu systému, nie z odpovedí AI asistentov.

AI nástroje boli teda využívané predovšetkým ako konzultačná a jazyková pomôcka, nie ako samostatný autor výsledného riešenia.

\section*{Záver}

Cieľom projektu bolo preskúmať možnosti odpovedania na otázky nad tabuľkovými dátami s využitím datasetu DataBench a navrhnúť reprodukovateľné riešenie, ktoré umožňuje automaticky vytvoriť odpoveď na základe zadanej otázky a príslušnej tabuľky. V rámci práce bol analyzovaný dataset DataBench, jeho štruktúra, typy otázok a požiadavky na vyhodnocovanie výsledkov.

V teoretickej časti boli preskúmané viaceré možnosti riešenia úlohy, vrátane špecializovaných modelov pre tabuľkové otázky, prístupu Text-to-SQL a generovania Python kódu. Na základe porovnania jednotlivých prístupov bol ako hlavný baseline zvolený prístup založený na generovaní Python kódu. Jeho výhodou je možnosť priamej práce s tabuľkovými dátami, riešenie rôznych typov operácií a zároveň možnosť zaznamenať a kontrolovať medzikrok medzi otázkou a výslednou odpoveďou.

V praktickej časti bol implementovaný pipeline, ktorý načíta otázku a príslušnú tabuľku, vytvorí prompt pre jazykový model, vygeneruje Python kód a následne ho vykoná v izolovanom prostredí. Výsledok vykonaného kódu je následne spracovaný do požadovaného formátu a môže byť vyhodnotený pomocou metrík definovaných pre dataset DataBench. Súčasťou riešenia je aj spracovanie chýb pri generovaní alebo vykonávaní kódu.

Navrhnuté riešenie vytvára základ pre ďalšie experimentovanie s rôznymi jazykovými modelmi, promptami a spôsobmi generovania odpovedí. V ďalšej fáze projektu bude možné porovnať výsledky hlavného baseline s priamym generovaním odpovedí a analyzovať najčastejšie typy chýb. Získané výsledky môžu zároveň ukázať, pri ktorých typoch otázok je generovanie kódu vhodnejšie a kde sú potrebné ďalšie úpravy alebo alternatívne prístupy.

% ZOZNAM POUŽITEJ LITERATÚRY


\begin{thebibliography}{9}

\bibitem{semeval}
SemEval-2025 Task 8: Question Answering over Tabular Data.
\newblock \url{https://aclanthology.org/2025.semeval-1.324/}

\bibitem{hfdatabench}
Hugging Face: \texttt{cardiffnlp/databench}.
\newblock \url{https://huggingface.co/datasets/cardiffnlp/databench}

\bibitem{databencheval}
\texttt{databench\_eval}: Evaluation framework for DataBench.
\newblock \url{https://github.com/jorses/databench_eval}

\bibitem{databench2024}
Osés Grijalba, J. et al. (2024).
Question Answering over Tabular Data with DataBench:
A Large-Scale Empirical Evaluation of LLMs.
In \emph{LREC-COLING 2024}.
\newblock \url{https://aclanthology.org/2024.lrec-main.1179/}

\bibitem{tapas}
Herzig, J. et al. (2020).
TaPas: Weakly Supervised Table Parsing via Pre-training.
In \emph{ACL 2020}.
\newblock \url{https://aclanthology.org/2020.acl-main.398/}

\bibitem{spider}
Yu, T. et al. (2018).
Spider: A Large-Scale Human-Labeled Dataset for Complex
and Cross-Domain Semantic Parsing and Text-to-SQL Task.
In \emph{EMNLP 2018}.
\newblock \url{https://aclanthology.org/D18-1425/}

\bibitem{tabularqa}
Evangelatos, A. et al.
Question Answering over Tabular Data with LLM Code Generation.
GitHub repository: \texttt{tabularqa}.
\newblock \url{https://github.com/andrewevag/tabularqa}

\bibitem{oseibrefo}
Osei-Brefo, E., Liang, H. (2025).
OseiBrefo-Liang at SemEval-2025 Task 8:
A Multi-Agent LLM Code Generation Approach for Answering Tabular Questions.
In \emph{Proceedings of the 19th International Workshop on Semantic Evaluation (SemEval-2025)}.
\newblock \url{https://aclanthology.org/2025.semeval-1.49/}

\end{thebibliography}

\end{document}
```